\subsection{Lower Bounds} 
Let's try to factor the polynomial
\[
3x^3+4x^2+13x+4
\]
\begin{itemize}
\item This polynomial has \makeblank[.5in] sign changes, so it has \makeblank[.5in] positive real roots!
\item Plugging in $-x$ for $x$ gives us \makeblank
\item This has \makeblank[.5in] sign changes, so our original polynomial has \makeblank[.5in] or \makeblank[.5in] negative real roots.
\end{itemize}
Its possible rational roots are:
\begin{lpic}{}{3}
\foreach \y/\val in {0/1,1/3}
  \regtextleft{1}{-2-\y}{$q=\pm\val$};
\foreach \x/\val in {0/1,1/2,2/4}
  \regtext{3+\value{cols}/3-1/2+\x*\value{cols}*2/3-\x}{-1}{$p=\pm\val$};
\end{lpic}
We might decide to try all the negative integers first.
\begin{cpic}{}{7}
	\foreach \a/\b in {-10/0,0/0,-10/-4} {
		  \synthdiv{\a}{\b}{4};
			\foreach \x/\lbl in {0/3,1/4,2/13,3/4}
				\regtextright{\a+2*\x+3}{\b-1}{$\lbl$};
			\regtextright{\a+3}{\b-3}{$3$};
	}
	\foreach \x/\y/\lbl in {-6/-2/ -3,-4/-2/ -1,-2/-2/ -12,-6/-3/ 1,-4/-3/ 12,-2/-3/  -8,
	                         4/-2/ -6, 6/-2/  4, 8/-2/ -34, 4/-3/-2, 6/-3/ 17, 8/-3/ -30,
													-6/-6/-12,-4/-6/ 32,-2/-6/-180,-6/-7/-8,-4/-7/ 45,-2/-7/-176,
												  -10/-1/-1,-10/-5/-4,0/-1/-2}
			\regtextright{\x+1}{\y}{$\lbl$};
	\draw[dashed] (0,0) -- (0,-7);
\end{cpic}
We might have worried that $-1$ was an upper bound, but it's not, for two reasons:
\begin{itemize}
\item Our rule for upper bounds only works for \makeblank[1in] numbers
\item The \makeblank[1in] is negative
\end{itemize}
However, there is another rule that helps us here.
\begin{theorem}
If we have a negative number $k$, and when performing synthetic division by $(x-k)$, the terms in the quotient and the remainder \emph{alternate} signs, $k$ is a \emph{lower bound} on the roots of the polynomial, that is, every root of the polynomial is greater than or equal to $k$.
\end{theorem}
In this example, \makeblank[.5in] was a lower bound, so we didn't need to try \makeblank[.5in]. The numbers \makeblank are still possibilities.